\subsection{Accelera Pipe Results}

The Accelera Pipe experiments compare three implementations under the same
train, validation, and test split. The first implementation is a manually
written scikit-learn loop over the candidate pipelines. The second uses
\texttt{sklearn.pipeline.Pipeline} to represent the same candidates. The third
uses Accelera Pipe with graph branches. Each candidate is selected using
validation accuracy, and the final selected path is evaluated on the test set.
This protocol checks whether the graph execution improves latency without
changing the selected model or the measured predictive accuracy.

\begin{table}[H]
\centering
\caption{Accelera Pipe latency and memory scaling}
\label{tab:pipe-scale}
\resizebox{\textwidth}{!}{
\begin{tabular}{rrrrrrrrrr}
\toprule
Samples & Train & Val & Test &
sklearn time & sklearn RSS &
Pipeline time & Pipeline RSS &
Accelera time & Accelera RSS \\
\midrule
1,000 & 600 & 200 & 200 &
0.06 & 1.62 & \bestcell 0.05 & \bestcell 0.16 & 0.34 & 19.32 \\
5,000 & 3,000 & 1,000 & 1,000 &
2.47 & 3.62 & 2.12 & \bestcell 1.09 & \bestcell 2.02 & 21.62 \\
10,000 & 6,000 & 2,000 & 2,000 &
2.48 & 11.47 & 2.48 & \bestcell 2.19 & \bestcell 1.79 & 7.25 \\
50,000 & 30,000 & 10,000 & 10,000 &
23.96 & 126.95 & 23.25 & \bestcell 9.75 & \bestcell 14.81 & 209.25 \\
100,000 & 60,000 & 20,000 & 20,000 &
77.79 & 139.70 & 77.10 & \bestcell 47.65 & \bestcell 67.04 & 300.93 \\
250,000 & 150,000 & 50,000 & 50,000 &
803.56 & 303.95 & 803.01 & \bestcell 28.52 & \bestcell 494.64 & 594.86 \\
\bottomrule
\end{tabular}}
\end{table}

Table~\ref{tab:pipe-scale} shows the main performance trend. At very small sizes, Accelera Pipe is
slower because the native graph setup, branch creation, and Python/C++ boundary
cost dominate the actual model computation. Once the dataset grows, those fixed
costs become small compared with fitting multiple candidate branches. At
50,000 samples, Accelera Pipe reduces runtime from 23.96 seconds to 14.81
seconds. At 250,000 samples, it reduces runtime from about 803 seconds to about
495 seconds, saving more than five minutes while selecting the same candidate.
The table also shows the current cost of graph-level parallelism: Accelera Pipe
uses more RSS memory because multiple branch states and intermediate arrays can
exist at the same time. The memory optimizations in the methodology section
target exactly this issue.

\begin{table}[H]
\centering
\caption{Selected candidate and accuracy across sample sizes}
\label{tab:pipe-accuracy}
\begin{tabular}{rlll}
\toprule
Samples & Selected candidate & Validation accuracy & Test accuracy \\
\midrule
1,000 & MinMaxScaler $\rightarrow$ SVC & 0.7900 & 0.8250 \\
5,000 & MinMaxScaler $\rightarrow$ SVC & 0.9160 & 0.9100 \\
10,000 & MinMaxScaler $\rightarrow$ SVC & 0.9385 & 0.9260 \\
50,000 & MinMaxScaler $\rightarrow$ SVC & 0.9565 & 0.9598 \\
100,000 & StandardScaler $\rightarrow$ SVC & 0.9675 & 0.9709 \\
250,000 & StandardScaler $\rightarrow$ SVC & 0.9774 & 0.9768 \\
\bottomrule
\end{tabular}
\end{table}

Table~\ref{tab:pipe-accuracy} confirms that the speedup does not come from changing the search
space. Accelera Pipe evaluates the same candidate pipelines and reports the
same validation and test behavior. The selected preprocessing strategy changes
from MinMax scaling to standard scaling only when the larger dataset makes
StandardScaler with SVC the best validation candidate. This is expected because
the candidate set and selection rule are fixed while the data distribution
becomes better represented at larger sizes.

\begin{table}[H]
\centering
\caption{Largest Accelera Pipe comparison}
\label{tab:pipe-largest}
\begin{tabular}{lrrrr}
\toprule
Implementation & Time (s) & RSS MB & VMS MB & Test accuracy \\
\midrule
sklearn manual loop & 803.56 & 303.95 & 315.41 & 0.9768 \\
sklearn Pipeline & 803.01 & \bestcell 28.52 & \bestcell 28.61 & 0.9768 \\
Accelera Pipe & \bestcell 494.64 & \warncell 594.86 & \warncell 787.47 & 0.9768 \\
\bottomrule
\end{tabular}
\end{table}

Table~\ref{tab:pipe-largest} isolates the largest run: 150,000 training samples, 50,000 validation
samples, and 50,000 test samples. Accelera Pipe keeps exactly the same selected
path, StandardScaler followed by SVC, and the same test accuracy. The runtime
drop is 308.91 seconds relative to the manual scikit-learn loop and 308.36
seconds relative to scikit-learn Pipeline. This result is important because it
demonstrates the value of graph scheduling on realistic branch-heavy workloads.
The memory columns also show the next engineering target: scikit-learn Pipeline
is extremely memory efficient, while Accelera Pipe trades memory for parallel
branch execution and native graph state.

\subsection{Code Parallelizer Results}

The OpenMP classifier was evaluated as a three-class loop classifier. The
labels are \texttt{none}, \texttt{parallel\_for}, and \texttt{reduction}. The
classification report shows that the model is balanced enough to guide the
rule layer, while the final pragma decision remains protected by dependency
checks.

\begin{table}[H]
\centering
\caption{OpenMP classifier report}
\label{tab:classifier-report}
\begin{tabular}{lrrrr}
\toprule
Class & Precision & Recall & F1-score & Support \\
\midrule
none & 0.82 & 0.85 & 0.84 & 3048 \\
parallel\_for & 0.86 & 0.81 & 0.83 & 2696 \\
reduction & 0.79 & 0.82 & 0.80 & 775 \\
\midrule
accuracy & & & 0.83 & 6519 \\
macro avg & 0.82 & 0.83 & 0.83 & 6519 \\
weighted avg & 0.83 & 0.83 & 0.83 & 6519 \\
\bottomrule
\end{tabular}
\end{table}

Table~\ref{tab:classifier-report} should be read as a guidance-quality result, not as a proof that every
predicted loop is safe. The classifier is responsible for identifying likely
parallelization opportunities. The rule layer is responsible for rejecting
unsafe or unsupported cases. This two-stage design is useful because source
code parallelization has correctness constraints that pure statistical
classification should not decide alone.

\begin{table}[H]
\centering
\caption{Hard parallelizer evaluation}
\label{tab:hard-eval}
\begin{tabular}{lrrrr}
\toprule
File & Baseline ms & Parallel ms & Speedup & Output match \\
\midrule
hard\_eval\_000.cpp & 1786.28 & \bestcell 195.47 & 9.138$\times$ & true \\
hard\_eval\_001.cpp & 63.98 & \bestcell 31.47 & 2.033$\times$ & true \\
hard\_eval\_002.py & 42.66 & \bestcell 7.67 & 5.563$\times$ & true \\
\bottomrule
\end{tabular}
\end{table}

Table~\ref{tab:hard-eval} evaluates harder programs from \modulepath{data/hard\_test\_files}.
The benchmark compiles and runs both the baseline and the generated
parallelized code, then compares their outputs. The first file obtains the
largest improvement because its loop body has enough work to amortize OpenMP
thread overhead. The second file still improves, but the smaller absolute
runtime limits the maximum speedup. The Python-origin file demonstrates the
full converter path: Python subset to C++, then loop extraction, then pragma
insertion. Across the three files, the evaluation extracted 11 loops, generated
8 pragmas, matched all 8 expected pragmas, reached average similarity 1.0, and
obtained an average speedup of 5.578$\times$.

\begin{table}[H]
\centering
\caption{Parallelizer and Accelera Pipe integration benchmark}
\label{tab:pipe-parallelizer-integration}
\begin{tabular}{lrr}
\toprule
Mode & Samples & Time (s) \\
\midrule
Python custom preprocessing function & 500,000 & 6.83 \\
End-to-end optimized run, first measured process & 500,000 & 5.88 \\
End-to-end optimized run, warmed method and module cache & 500,000 & \bestcell 0.08 \\
\bottomrule
\end{tabular}
\end{table}

Table~\ref{tab:pipe-parallelizer-integration} reports the direct
\modulepath{examples/parallel\_accpipe.py} run after normalizing the example
execution directory to the repository root. The Python implementation takes
6.83 seconds on 500,000 samples. The first optimized process in the measurement
sequence takes 5.88 seconds because it still pays process-local setup and cache
lookup costs. The immediate repeated run takes 0.08 seconds after the
Python-method cache and compiled-module cache are both warm. The validation in
the example confirmed that the optimized output matched the Python output. This
result shows why the integration needs two cache layers: one layer avoids
repeating Python-to-C++ conversion and OpenMP insertion, while the lower
compiled-module cache avoids recompiling the pybind extension.

\begin{table}[H]
\centering
\caption{Direct example-script verification runs}
\label{tab:direct-example-runs}
\resizebox{\textwidth}{!}{
\begin{tabular}{lrrr}
\toprule
Example & Baseline time (s) & Optimized time (s) & Validation \\
\midrule
\modulepath{accpipe\_demo.py} & 2.48 & \bestcell 1.79 & same test accuracy, 0.9260 \\
\modulepath{parallel\_accpipe.py}, run 1 & 6.83 & \bestcell 5.88 & outputs match \\
\modulepath{parallel\_accpipe.py}, run 2 & 6.83 & \bestcell 0.08 & outputs match \\
\modulepath{code\_parallelizer\_demo.py}, run 1 & 5.81 & \bestcell 0.014 & outputs match \\
\modulepath{code\_parallelizer\_demo.py}, run 2 & 6.05 & \bestcell 0.015 & outputs match \\
\modulepath{save\_load.py}, run 1 & 1.89 & 2.69 / 2.65 & same accuracy, 0.268 \\
\modulepath{save\_load.py}, run 2 & 2.53 & \bestcell 2.20 / 2.14 & same accuracy, 0.268 \\
\bottomrule
\end{tabular}}
\end{table}

Table~\ref{tab:direct-example-runs} summarizes the examples invoked by
\modulepath{examples/Makefile}, but each file was run directly rather than
through Make. Cache-sensitive files were repeated. The standalone
\modulepath{code\_parallelizer\_demo.py} result measures Python script execution
against the generated OpenMP executable for \modulepath{hard\_eval\_002.py}; the
C path also reported about 0.26--0.27 seconds of compilation time, outside the
listed executable runtime. The save/load example verifies persistence rather
than pure acceleration: both sklearn and Accelera load persisted state and
produce the same accuracy, while the second Accelera run is slightly faster
than the sklearn baseline.
